\section{Differentiable Manifolds}

\noindent\textbf{Exercise~\ref{ex:diffmfds-1}.}\label{sol:diffmfds-1}

We have $x^{-1}:\R \to \R$; $a\mapsto x^{-1}(a) = a$, and so $y\circ x^{-1}:\R\to\R$; $a\mapsto a^3$. This is $C^{\infty}(\R)$.

\bigskip
\noindent\textbf{Exercise~\ref{ex:diffmfds-2}.}\label{sol:diffmfds-2}

We have $x\circ y^{-1}: \R\to\R$; $a\mapsto \sqrt[3]{a}$, but this function is not even $C^1(\R\to\R)$ as the first derivative is $\frac{1}{3}a^{-2/3}$, which blows up as $a\to 0$. So no it is not a differentiable manifold.

\bigskip
\noindent\textbf{Exercise~\ref{ex:diffmfds-3}.}\label{sol:diffmfds-3}

The problem above came from the chart transition maps. Both chart maps themselves are $C^{\infty}(\R\to\R)$ and their domain is the whole manifold (namley $\R$) and so if we just remove one of the two charts we get a smooth manifold.

\bigskip
\noindent\textbf{Exercise~\ref{ex:diffmfds-4}.}\label{sol:diffmfds-4}

For $(U_n,x_n)$ to be charts, we need to show that $U_n\in \cO_{ssq}$ and that the $x_n$s are homeomorphisms (i.e. inevitable and continuous).

The soft square topology is rather self explanatory, it is the set of squares around the origin without the boundary. This is exactly the definition of the $U_n$s where the sides of consecutive $U_n$ increase by $2n$. So we have $U_n\in\cO_{ssq}$.

Next we need to show that the chart maps are homeomorphisms. It is clear that $x_n$ is continuous w.r.t. $\cO_{ssq}$ and $\cO_{\text{standard}}$. So we just need to check that $x_n^{-1}$ exists and is continuous. We have
\bse
    x_n^{-1}(a,b) = \big(n(a+b), n(a-b)\big),
\ese
which is again clearly continuous.

Therefore we know that the $(U_n,x_n)$ are indeed charts.

\bigskip
\noindent\textbf{Exercise~\ref{ex:diffmfds-5}.}\label{sol:diffmfds-5}

The chart transition maps are given by
\bse
    \begin{split}
        x_m \circ x_n^{-1} : x_n(U_n\cap U_m) & \to x_m(U_n\cap U_m) \\
        (a,b) & \mapsto \bigg(\frac{na}{m}, \frac{nb}{m}\bigg) = \frac{n}{m}(a,b).
    \end{split}
\ese
This tells us that $x_m\circ x_n^{-1} = \frac{n}{m}\b1_{\R^2}$, which together with $m\neq 0$ tells us $k=\infty$.

\bigskip
\noindent\textbf{Exercise~\ref{ex:diffmfds-6}.}\label{sol:diffmfds-6}

There are many, but we could take the map $\widetilde{x}:U_5 \to \widetilde{U_5}\ss\R^2$; $(a,b) \mapsto (a,b)$, i.e. it is the identity map restricted to $U_5$. It is clear that this will be $C^{\infty}$ compatible with any overlapping charts and so it lies in the atlas.
